\section{Matching normal forms}\label{sec:normalforms}

The engine of both sufficiency proofs is a reduction of $A$ to a sparse
normal form by invertible triangular transformations.  The normal forms are
"matrix matchings": in the unsymmetric case a matching between row indices
and column indices, in the symmetric case a matching of the index set with
itself.

\begin{definition}[matching matrix]\label{def:matching}
  A matrix $\Pi \in \F^{n \by n}$ is a \emph{matching matrix} if every row
  and every column of $\Pi$ contains at most one nonzero entry.
  Equivalently,
  \[
    \Pi = \sum_{(i,j) \in \M} w_{ij} \, e_i e_j^T,
  \]
  where the \emph{matching} $\M = \M(\Pi) \subseteq \{1,\ldots,n\}^2$ has
  pairwise distinct first coordinates and pairwise distinct second
  coordinates, and the weights $w_{ij} \in \F$ are nonzero.  We call
  $(i,j) \in \M$ a \emph{matched pair}, a row not appearing as a first
  coordinate a \emph{zero row}, and a column not appearing as a second
  coordinate a \emph{zero column}.  When all weights equal $1$, $\Pi$ is a
  partial permutation matrix.
\end{definition}

\begin{definition}[symmetric matching matrix]\label{def:sym-matching}
  A symmetric matrix $P \in \F^{n \by n}$ is a \emph{symmetric matching
  matrix} if the index set $\{1, \ldots, n\}$ can be partitioned into
  \begin{itemize}
    \item \emph{isolated zero} indices, whose rows and columns vanish;
    \item \emph{singletons} $t$ with $P[t,t] = \gamma_t \neq 0$ and no other
      nonzero entry in row or column $t$;
    \item \emph{pairs} $(i,j)$ with $i < j$,
      \[
        P[i,j] = P[j,i] = \alpha \neq 0,
        \qquad
        P[i,i] = 0,
        \qquad
        P[j,j] = \beta \in \F ,
      \]
      and no other nonzero entries in rows and columns $i, j$.
  \end{itemize}
  Each pair block
  $\bigl(\begin{smallmatrix} 0 & \alpha\\ \alpha & \beta
  \end{smallmatrix}\bigr)$ is nonsingular (determinant $-\alpha^2 \neq 0$).
  Note the asymmetry within a pair: the \emph{earlier} index $i$ carries a
  zero diagonal entry, while the later index $j$ may carry an arbitrary
  diagonal weight $\beta$.
\end{definition}

The nonzero pattern of a symmetric matching matrix is not literally a
matching when $\beta \neq 0$ (row $j$ then carries two nonzero entries), but
its \emph{combinatorial content} is the matching that pairs $i$ with $j$ and
matches each singleton with itself; \Cref{prop:canonical} below makes this
precise.

\subsection{Reduction of an arbitrary matrix}

\begin{lemma}[triangular equivalence to a matching matrix]\label{lem:reduction-lu}
  For every $A \in \F^{n \by n}$ there exist an invertible lower triangular
  $M$, an invertible upper triangular $N$, and a matching matrix $\Pi$ such
  that
  \[
    A = M \, \Pi \, N .
  \]
  The construction below computes $M$, $\Pi$, $N$ in $O(n^3)$ field
  operations, and $M$, $N$ can be taken unit triangular.
\end{lemma}

\begin{proof}
  We induct on $k$: we construct invertible unit lower (resp.\ upper)
  triangular $M^{(k)}, N^{(k)} \in \F^{k \by k}$ and matching matrices
  $\Pi^{(k)}$ with $A[1:k,1:k] = M^{(k)} \Pi^{(k)} N^{(k)}$.  For $k = 1$
  take $M^{(1)} = N^{(1)} = (1)$ and $\Pi^{(1)} = (a_{11})$, whose matching
  is $\{(1,1)\}$ if $a_{11} \neq 0$ and empty otherwise.

  Assume the claim for $k - 1$, write $M = M^{(k-1)}$, $\Pi = \Pi^{(k-1)}$,
  $N = N^{(k-1)}$, $\M = \M(\Pi)$, and border
  \[
    A[1:k,1:k] =
    \begin{pmatrix}
      A[1:k-1,1:k-1] & a\\
      c^T & b
    \end{pmatrix},
    \qquad a, c \in \F^{k-1}, \; b \in \F .
  \]
  Setting $\widehat a := M^{-1} a$ and $\widehat c^{\,T} := c^T N^{-1}$,
  \begin{equation}\label{eq:border-lu}
    A[1:k,1:k] =
    \begin{pmatrix}
      M & 0\\
      0 & 1
    \end{pmatrix}
    \begin{pmatrix}
      \Pi & \widehat a\\
      \widehat c^{\,T} & b
    \end{pmatrix}
    \begin{pmatrix}
      N & 0\\
      0 & 1
    \end{pmatrix},
  \end{equation}
  and it remains to reduce the bordered middle matrix by further unit
  triangular factors.  The reduction has two stages.

  \emph{Stage 1: strip the matched part of the border.}
  Define $x, y \in \F^{k-1}$ supported on the matched columns and rows of
  $\Pi$:
  \[
    x_j := \widehat a_i / w_{ij},
    \qquad
    y_i := \widehat c_j / w_{ij},
    \qquad (i,j) \in \M,
  \]
  with all other components zero.  Then $\Pi x$ agrees with $\widehat a$ on
  every matched row, and $y^T \Pi$ agrees with $\widehat c^{\,T}$ on every
  matched column, so
  \begin{gather*}
    e := \widehat a - \Pi x
    \quad \text{is supported on the zero rows of } \Pi,\\
    g^T := \widehat c^{\,T} - y^T \Pi
    \quad \text{is supported on the zero columns of } \Pi .
  \end{gather*}
  Moreover $y^T e = 0$ and $g^T x = 0$, because in each product the supports
  are disjoint (matched versus zero indices).  Setting
  $f := b - y^T \Pi \, x$, a direct block multiplication gives
  \begin{equation}\label{eq:peel-lu}
    \begin{pmatrix}
      \Pi & \widehat a\\
      \widehat c^{\,T} & b
    \end{pmatrix}
    =
    \begin{pmatrix}
      I & 0\\
      y^T & 1
    \end{pmatrix}
    \begin{pmatrix}
      \Pi & e\\
      g^T & f
    \end{pmatrix}
    \begin{pmatrix}
      I & x\\
      0 & 1
    \end{pmatrix}:
  \end{equation}
  indeed the product equals
  $\bigl(\begin{smallmatrix} \Pi & \Pi x + e\\ y^T \Pi + g^T & y^T \Pi x +
  y^T e + g^T x + f \end{smallmatrix}\bigr)$, and the corner reduces to
  $y^T \Pi x + f = b$.

  \emph{Stage 2: absorb the residual border.}
  The middle matrix of \cref{eq:peel-lu} now has its border supported on
  previously zero rows and columns of $\Pi$; we convert it into at most two
  new matched pairs.

  \emph{(2a) Residual column.}
  If $e \neq 0$, let $i^\star$ be the index of its \emph{first} nonzero
  component and $\alpha := e_{i^\star}$; by Stage 1, $i^\star$ is a zero row
  of $\Pi$.  Let
  \[
    R := I + (\alpha^{-1} e - u_{i^\star}) \, u_{i^\star}^T \in
    \F^{(k-1) \by (k-1)},
  \]
  where $u_{i^\star}$ is the $i^\star$th standard basis vector.  Since the
  first nonzero of $e$ sits at $i^\star$, the vector
  $\alpha^{-1}e - u_{i^\star}$ vanishes at and above position $i^\star$, so
  $R$ is unit lower triangular.  Because row $i^\star$ of $\Pi$ is zero,
  $R \, \Pi = \Pi$; and $R \, u_{i^\star} = \alpha^{-1} e$.  Hence
  \[
    \begin{pmatrix}
      \Pi & e\\
      g^T & f
    \end{pmatrix}
    =
    \begin{pmatrix}
      R & 0\\
      0 & 1
    \end{pmatrix}
    \begin{pmatrix}
      \Pi & \alpha u_{i^\star}\\
      g^T & f
    \end{pmatrix}.
  \]
  If $e = 0$, skip this step.

  \emph{(2b) Residual row.}
  Symmetrically, if $g \neq 0$, let $j^\star$ be the first nonzero index of
  $g$ and $\gamma := g_{j^\star}$; $j^\star$ is a zero column of $\Pi$.  With
  the unit upper triangular
  $S := I + u_{j^\star} (\gamma^{-1} g - u_{j^\star})^T$ we have
  $\Pi S = \Pi$ and $u_{j^\star}^T S = \gamma^{-1} g^T$, so
  \[
    \begin{pmatrix}
      \Pi & v\\
      g^T & f
    \end{pmatrix}
    =
    \begin{pmatrix}
      \Pi & v\\
      \gamma u_{j^\star}^T & f
    \end{pmatrix}
    \begin{pmatrix}
      S & 0\\
      0 & 1
    \end{pmatrix},
    \qquad v \in \{\alpha u_{i^\star}, 0\},
  \]
  where we used that $S$ acts on the columns $1, \ldots, k-1$ only, leaving
  the border column $v$ untouched.

  \emph{(2c) Corner.}
  Let $X$ denote the current middle matrix, with border column
  $v = \alpha u_{i^\star}$ or $0$, border row $\gamma u_{j^\star}^T$ or $0$,
  and corner $f$.  Three cases:
  \begin{itemize}
    \item If $e \neq 0$: row $i^\star$ of $X$ equals $(0, \ldots, 0 \mid
      \alpha)$, so with the unit lower triangular
      $T := I_k + (f/\alpha) \, u_k u_{i^\star}^T$ we get $X = T X_0$, where
      $X_0$ is $X$ with corner $0$.
    \item Else if $g \neq 0$: column $j^\star$ of $X$ equals
      $\gamma \, u_k$, so with the unit upper triangular
      $T' := I_k + (f/\gamma) \, u_{j^\star} u_k^T$ we get $X = X_0 T'$.
    \item Else $X = \Pi \oplus (f)$.
  \end{itemize}

  Collecting the factors of Stages 1--2 into the outer factors of
  \cref{eq:border-lu} yields
  $A[1:k,1:k] = M^{(k)} \, \Pi^{(k)} \, N^{(k)}$ with unit triangular
  invertible $M^{(k)}, N^{(k)}$ and
  \[
    \M(\Pi^{(k)}) = \M(\Pi) \;\cup\;
    \begin{cases}
      \{(i^\star, k)\} & \text{if } e \neq 0 \text{ (weight } \alpha)\\
      \emptyset & \text{else}
    \end{cases}
    \;\cup\;
    \begin{cases}
      \{(k, j^\star)\} & \text{if } g \neq 0 \text{ (weight } \gamma)\\
      \{(k,k)\} & \text{if } e = 0,\ g = 0,\ f \neq 0\\
      \emptyset & \text{else.}
    \end{cases}
  \]
  This is again a matching: the new pairs occupy row $i^\star$ (previously a
  zero row), column $j^\star$ (previously a zero column), and the new index
  $k$, whose row is used by at most one of $(k, j^\star)$, $(k,k)$ and whose
  column by at most one of $(i^\star, k)$, $(k,k)$.  Taking $k = n$
  completes the induction.  Each border step costs $O(k^2)$ operations (two
  triangular solves and rank-one updates of $M$ and $N$), for $O(n^3)$ in
  total.  Only additions, multiplications, and divisions by the nonzero
  scalars $w_{ij}, \alpha, \gamma$ are used, so the construction is valid
  over every field.
\end{proof}

\begin{remark}[normalizing the weights]\label{rem:normalize-weights}
  Writing $W$ for the diagonal matrix with $W[j,j] = w_{ij}$ when column $j$
  is matched by $(i,j) \in \M$ and $W[j,j] = 1$ otherwise, we have
  $\Pi = \Pi_0 W$ with $\Pi_0$ a partial permutation matrix; absorbing $W$
  into $N$ shows that the middle factor in \Cref{lem:reduction-lu} can be
  taken to be a $0$--$1$ partial permutation matrix.  For invertible $A$
  there are no zero rows or columns, $\Pi_0$ is a permutation, and
  $A = M \Pi_0 N'$ is the Bruhat (LPU)
  decomposition~\cite{elsner_bruhat_1979,strang_lpu_2012}; for singular $A$
  it is the generalized Bruhat (LEU)
  decomposition~\cite{malaschonok_2010,dumas_pernet_sultan_2017}.
\end{remark}

\subsection{Reduction of a symmetric matrix}

Under congruence the two outer factors are transposes of one another, and
weights can no longer be normalized away (only multiplied by squares); this
is why the symmetric normal form retains the weights $\alpha$ and the
diagonal fill $\beta$.

\begin{lemma}[triangular congruence to a symmetric matching matrix]\label{lem:reduction-sym}
  For every symmetric $A \in \F^{n \by n}$ there exist an invertible lower
  triangular $M$ and a symmetric matching matrix $P$ such that
  \[
    A = M \, P \, M^T .
  \]
  The construction below computes $M$, $P$ in $O(n^3)$ field operations,
  and $M$ can be taken unit lower triangular.
\end{lemma}

\begin{proof}
  As before we induct on $k$, with base case
  $A[1:1,1:1] = (1)(a_{11})(1)$: the index $1$ is a singleton if
  $a_{11} \neq 0$ and an isolated zero otherwise.  Assume
  $A[1:k-1,1:k-1] = M P M^T$ with $M$ unit lower triangular and $P$ a
  symmetric matching matrix, and border
  \[
    A[1:k,1:k] =
    \begin{pmatrix}
      A[1:k-1,1:k-1] & a\\
      a^T & b
    \end{pmatrix},
    \qquad
    c := M^{-1} a,
  \]
  so that
  \begin{equation}\label{eq:border-sym}
    A[1:k,1:k] =
    \begin{pmatrix}
      M & 0\\
      0 & 1
    \end{pmatrix}
    \begin{pmatrix}
      P & c\\
      c^T & b
    \end{pmatrix}
    \begin{pmatrix}
      M^T & 0\\
      0 & 1
    \end{pmatrix}.
  \end{equation}

  \emph{Stage 1: strip the matched part of the border.}
  $P$ is invertible on each of its nonzero blocks, so we may choose
  $\xi \in \F^{k-1}$, supported on the singleton and pair indices, with
  $P \xi$ equal to $c$ there: for a singleton $t$ set
  $\xi_t = c_t / \gamma_t$; for a pair $(i,j)$ with weights
  $(\alpha, \beta)$ solve the $2 \by 2$ system, i.e.
  $\xi_j = c_i / \alpha$ and $\xi_i = (c_j - \beta \xi_j)/\alpha$.  Then
  \[
    e := c - P\xi
  \]
  is supported on the isolated zero indices of $P$, and $\xi^T e = 0$
  (disjoint supports).  With $f := b - \xi^T P \xi$, a direct block
  multiplication (using $\xi^T e = e^T \xi = 0$ in the corner) gives
  \begin{equation}\label{eq:peel-sym}
    \begin{pmatrix}
      P & c\\
      c^T & b
    \end{pmatrix}
    =
    \begin{pmatrix}
      I & 0\\
      \xi^T & 1
    \end{pmatrix}
    \begin{pmatrix}
      P & e\\
      e^T & f
    \end{pmatrix}
    \begin{pmatrix}
      I & \xi\\
      0 & 1
    \end{pmatrix}.
  \end{equation}

  \emph{Stage 2: absorb the residual border.}
  If $e = 0$, the middle matrix is $P \oplus (f)$: append the singleton $k$
  with weight $\gamma_k = f$ if $f \neq 0$, or the isolated zero $k$ if
  $f = 0$.  If $e \neq 0$, let $i^\star$ be its first nonzero index and
  $\alpha := e_{i^\star} \neq 0$; since $e$ is supported on isolated
  zeros, row and column $i^\star$ of $P$ vanish.  The unit lower triangular
  $R := I + (\alpha^{-1} e - u_{i^\star}) u_{i^\star}^T$ satisfies
  $R P R^T = P$ and $R u_{i^\star} = \alpha^{-1} e$, whence
  \[
    \begin{pmatrix}
      P & e\\
      e^T & f
    \end{pmatrix}
    =
    \begin{pmatrix}
      R & 0\\
      0 & 1
    \end{pmatrix}
    \begin{pmatrix}
      P & \alpha u_{i^\star}\\
      \alpha u_{i^\star}^T & f
    \end{pmatrix}
    \begin{pmatrix}
      R^T & 0\\
      0 & 1
    \end{pmatrix}.
  \]
  The middle matrix is the symmetric matching matrix obtained from $P$ by
  adding the pair $(i^\star, k)$ with weights $\alpha$ and $\beta = f$: the
  earlier index $i^\star$ (previously isolated) has zero diagonal entry, as
  required by \Cref{def:sym-matching}.  Absorbing the congruence factors of
  \cref{eq:peel-sym} and Stage~2 into the outer factors of
  \cref{eq:border-sym} completes the induction.  The cost analysis and
  field-generality are as in \Cref{lem:reduction-lu}; in particular no
  division by $2$ occurs, so the construction is valid in characteristic
  $2$.
\end{proof}

\begin{remark}[relation to known canonical forms]\label{rem:known-forms}
  Symmetric elimination without pivoting was studied by Van Dooren and
  Strang~\cite{vandooren_strang_2014}, who factor a real symmetric matrix
  as $A = LPL^T$ with $P$ built from diagonal entries $0, \pm 1$ and
  symmetric exchange blocks; \Cref{lem:reduction-sym} is a field-general,
  weighted version of that reduction.  From the canonical-form standpoint,
  Szechtman~\cite{szechtman_2007} classified symmetric matrices under
  congruence by triangular groups: for $\operatorname{char} \F \neq 2$ the
  canonical forms under unit-triangular congruence are weighted symmetric
  sub-permutation matrices, over any field, while for fields of
  characteristic $2$ (assuming every element is a square) the list of
  canonical forms necessarily includes $2 \by 2$ blocks with a nonzero
  corner --- precisely the $\beta$-entries retained in
  \Cref{def:sym-matching}.  The proof given above is included because it is
  short, uniform in the field and the characteristic, and is literally the
  algorithm implemented in \Cref{sec:validation}.
\end{remark}

\begin{remark}[signed matching matrices over $\mathbb{R}$]\label{rem:signed}
  Over $\mathbb{R}$ (more generally, over any field of characteristic
  $\neq 2$ in which every element or its negative is a square) the
  symmetric normal form can be normalized further by a diagonal-plus-shear
  congruence, applied blockwise.  For a pair block, with $i < j$,
  \[
    \begin{pmatrix}
      1 & 0\\
      -\beta/(2\alpha) & 1
    \end{pmatrix}
    \begin{pmatrix}
      0 & \alpha\\
      \alpha & \beta
    \end{pmatrix}
    \begin{pmatrix}
      1 & -\beta/(2\alpha)\\
      0 & 1
    \end{pmatrix}
    =
    \begin{pmatrix}
      0 & \alpha\\
      \alpha & 0
    \end{pmatrix},
  \]
  where the shear adds a multiple of row/column $i$ to row/column $j > i$
  and is therefore lower triangular; scaling index $j$ by $\alpha^{-1}$
  turns the block into
  $\bigl(\begin{smallmatrix} 0 & 1\\ 1 & 0 \end{smallmatrix}\bigr)$, and
  scaling a singleton $t$ by $|\gamma_t|^{-1/2}$ turns it into
  $\sgn(\gamma_t)$.  Every real symmetric matrix is thus congruent, by an
  invertible lower triangular congruence, to a \emph{signed matching
  matrix}: a symmetric matrix with entries in $\{0, \pm 1\}$, at most one
  nonzero per row and column, diagonal entries in $\{0, \pm 1\}$, and
  off-diagonal exchange blocks
  $\bigl(\begin{smallmatrix} 0 & 1\\ 1 & 0 \end{smallmatrix}\bigr)$.  The
  division by $2$ makes this normalization unavailable in characteristic
  $2$, which is why \Cref{lem:reduction-sym} retains $\beta$; over
  $GF(2)$, e.g., $\bigl(\begin{smallmatrix} 0 & 1\\ 1 & 1
  \end{smallmatrix}\bigr)$ is not congruent to
  $\bigl(\begin{smallmatrix} 0 & 1\\ 1 & 0 \end{smallmatrix}\bigr)$: the
  former is nonalternating, the latter alternating, and congruence
  preserves this distinction.
\end{remark}

\subsection{The normal form is canonical}

The matchings produced by
\Cref{lem:reduction-lu,lem:reduction-sym} are not artifacts of the
particular construction: they are invariants of $A$, computable from its
rank profile.  Part~(1) below coincides with the \emph{rank profile matrix}
of Dumas, Pernet, and Sultan~\cite{dumas_pernet_sultan_2017}; part~(2) is
the congruence analogue (cf.~\cite{bagno_cherniavsky_2012,szechtman_2007}).
For $1 \le i, j \le n$ let
\[
  r_{ij}(A) := \rk\bigl(A[1:i,1:j]\bigr),
  \qquad
  \Delta_{ij}(A) := r_{ij} - r_{i-1,j} - r_{i,j-1} + r_{i-1,j-1},
\]
with the convention $r_{0j} = r_{i0} = 0$.  We call $(i,j)$ a
\emph{rank jump} of $A$ if $\Delta_{ij}(A) = 1$; that
$\Delta_{ij} \in \{0,1\}$ always holds follows from
\Cref{lem:reduction-lu} together with part~(1) below (or
see~\cite{dumas_pernet_sultan_2017}).

\begin{proposition}[canonicity]\label{prop:canonical}
  \begin{enumerate}
    \item If $A = M \Pi N$ as in \Cref{lem:reduction-lu}, then
      \[
        r_{ij}(A) = \#\{(i',j') \in \M(\Pi) : i' \le i,\ j' \le j\}
        \quad \text{for all } i, j,
      \]
      and consequently $\M(\Pi)$ is exactly the set of rank jumps of $A$.
      In particular the matching --- and hence the zero rows and zero
      columns of $\Pi$ --- is uniquely determined by $A$ (the weights are
      not: diagonal scalings can be absorbed into $M$ or $N$).
    \item If $A = A^T = M P M^T$ as in \Cref{lem:reduction-sym}, then the
      singletons of $P$ are the diagonal rank jumps $\{t :
      \Delta_{tt}(A) = 1\}$ and the pairs of $P$ are the unordered
      off-diagonal rank jumps $\{\{i,j\} : i \neq j,\ \Delta_{ij}(A) = 1\}$;
      the partition of \Cref{def:sym-matching} is uniquely determined by
      $A$.
  \end{enumerate}
\end{proposition}

\begin{proof}
  (1) By \cref{eq:tri-slices},
  $A[1:i,1:j] = M[1:i,:] \, \Pi \, N[:,1:j] = M_i \, \Pi[1:i,1:j] \, N_j$
  with $M_i, N_j$ invertible, so $r_{ij}(A) = r_{ij}(\Pi)$ by
  \Cref{lem:rank-invariance}.  The nonzero columns of $\Pi[1:i,1:j]$ are
  the vectors $w_{i'j'} e_{i'}$ for the matched pairs with
  $i' \le i,\ j' \le j$; they occupy distinct rows, hence are linearly
  independent, giving the rank formula.  The function
  $(i,j) \mapsto \#\{(i',j') \in \M : i' \le i, j' \le j\}$ is a sum of
  corner indicators $\mathbf{1}_{\{i \ge i',\, j \ge j'\}}$, and
  $\Delta_{ij}(\mathbf{1}_{\{i \ge a,\, j \ge b\}}) = \mathbf{1}_{\{(i,j) =
  (a,b)\}}$, so the rank jumps of $A$ are exactly the matched pairs.

  (2) Congruence is the special case $N = M^T$ of a triangular
  equivalence, so as in (1), $r_{ij}(A) = r_{ij}(P)$ for all $i,j$.  The
  nonzero blocks of $P$ occupy pairwise disjoint row sets and column sets,
  so any rectangle $P[1:i,1:j]$ decomposes, after row and column
  permutations, into a direct sum of the blocks' restrictions, and its rank
  is the sum of the blocks' contributions.  A singleton $t$ contributes
  $\mathbf{1}_{\{i \ge t,\, j \ge t\}}$.  A pair $(i_0, j_0)$, $i_0 < j_0$,
  restricts to a matrix of rank $2$ when both its rows and columns are
  present, i.e.\ on the region $\{i \ge j_0,\ j \ge j_0\}$ (the two entries
  $\alpha$ are then present and the block has determinant
  $-\alpha^2 \neq 0$, regardless of $\beta$); it restricts to a matrix of
  rank $\ge 1$ exactly when at least one entry $\alpha$ is present, i.e.\ on
  the union $\{i \ge i_0,\ j \ge j_0\} \cup \{i \ge j_0,\ j \ge i_0\}$,
  whose intersection is $\{i \ge j_0,\ j \ge j_0\}$.  Writing the rank
  contribution as
  $\mathbf{1}_{\{\text{rank} \ge 1\}} + \mathbf{1}_{\{\text{rank} = 2\}}$
  and using inclusion--exclusion for the union, the contribution of the pair
  equals
  \[
    \Bigl(
    \mathbf{1}_{\{i \ge i_0,\, j \ge j_0\}}
    + \mathbf{1}_{\{i \ge j_0,\, j \ge i_0\}}
    - \mathbf{1}_{\{i \ge j_0,\, j \ge j_0\}}
    \Bigr)
    + \mathbf{1}_{\{i \ge j_0,\, j \ge j_0\}}
    =
    \mathbf{1}_{\{i \ge i_0,\, j \ge j_0\}}
    + \mathbf{1}_{\{i \ge j_0,\, j \ge i_0\}} .
  \]
  Taking mixed
  differences, the rank jumps contributed are exactly $(t,t)$ for
  singletons and $(i_0, j_0), (j_0, i_0)$ for pairs --- in particular
  $\beta$ produces no jump at $(j_0, j_0)$.  Since blocks are disjoint, the
  jump sets are disjoint, and the claim follows.
\end{proof}

\begin{remark}[compatibility of the two normal forms]\label{rem:compatibility}
  Let $A$ be symmetric.  By \Cref{prop:canonical}, the rank-jump set of $A$
  is symmetric ($\Delta_{ij} = \Delta_{ji}$), and the unsymmetric normal
  form $\Pi$ of $A$ has matching
  \[
    \M(\Pi) = \{(t,t) : t \text{ singleton of } P\}
    \;\cup\;
    \{(i,j), (j,i) : (i,j) \text{ pair of } P\} .
  \]
  Thus each $2 \by 2$ pair block of the congruence normal form corresponds
  to the \emph{two} matched pairs $(i,j)$ and $(j,i)$ of the equivalence
  normal form.  This dictionary is what makes
  \Cref{cor:sym-lu-ldlt} structurally transparent, and we will see it again
  in the combinatorial conditions of \Cref{sec:combinatorics}.
\end{remark}
